\documentclass[11pt,a4paper]{article}

%% PACHETE MATEMATICE
\usepackage{amsmath,amssymb,amsfonts}
\usepackage{amsthm}
\usepackage{mathpazo}
\usepackage{eulervm}
\usepackage{enumerate}
\usepackage{color}
\usepackage{graphicx}
\usepackage[all]{xy}
\usepackage{xcolor}
\usepackage{framed}
\usepackage[unicode]{hyperref}
\setcounter{MaxMatrixCols}{20}
\usepackage{multicol}
%\usepackage[romanian]{babel}


%% ASPECT
\linespread{1.05}
\usepackage[T1]{fontenc}
\usepackage[utf8x]{inputenc}
\usepackage[margin=2cm]{geometry}
% \usepackage[color]{showkeys}
\usepackage{anyfontsize}
\usepackage{titlesec}
\titleformat*{\section}{\large\bfseries}

\usepackage{fancyhdr}
\pagestyle{fancy}
\rhead{\small \color{gray} Notițe de Adrian Manea}
\lhead{\small \color{gray} BCD-Aero, 2017---2018, L. Costache \& M. Olteanu}


\usepackage[autostyle = false, style = german]{csquotes}
\usepackage[romanian]{babel}
\newcommand\qq{\enquote}

%% COMENZILE MELE
\renewcommand*{\proofname}{Dem.:}
\newcommand\bb{\mathbb}
\newcommand\dr{\mathrm}
\newcommand\kal{\mathcal}
\newcommand\fk{\mathfrak}
\newcommand\op{\oplus}
\newcommand\ot{\otimes}
\newcommand\ds{\displaystyle}
\newcommand\id{\indent}
\newcommand\nid{\noindent}
\newcommand\seq{\subseteq}
\newcommand\sneq{\subsetneq}
\newcommand\speq{\supseteq}
\newcommand\spneq{\supsetneq}
\newcommand\rar{\rightarrow}
\newcommand\Rar{\Rightarrow}
\newcommand\xrar{\xrightarrow}
\newcommand\lar{\leftarrow}
\newcommand\Lar{\Leftarrow}
\newcommand\xlar{\xleftarrow}
\newcommand\lrar{\leftrightarrow}
\newcommand\Lrar{\Leftrightarrow}
\newcommand\ideal{\trianglelefteq}
\newcommand\ai{\text{ a.\^{i}. }}
\newcommand\sk{(\textit{Schi\c{t}\u{a}}) }
\newcommand\rhup{\rightharpoonup}
\newcommand\rhdn{\rightharpoondown}
\newcommand\lhup{\leftharpoonup}
\newcommand\lhdn{\leftharpoondown}
\newcommand\vid{\emptyset}
\newcommand\wh{\widehat}
\newcommand\ol{\overline}
\newcommand\NN{\mathbb{N}}
\newcommand\RR{\mathbb{R}}
\newcommand\ZZ{\mathbb{Z}}
\newcommand\QQ{\mathbb{Q}}
\newcommand\CC{\mathbb{C}}

%% STILURILE MELE
\newtheoremstyle{dotless}{}{}{\itshape}{}{\bfseries}{:}{ }{}

\theoremstyle{dotless}
\newtheorem{theorem}{Teorem\u{a}}[section]
\newtheorem{proposition}{Propozi\c{t}ie}[section]
\newtheorem{lemma}{Lem\u{a}}[section]
\newtheorem{corollary}{Corolar}[section]
\newtheorem{exercise}{Exerci\c{t}iu}

\newtheoremstyle{dotlessdr}{}{}{}{}{\bfseries}{:}{ }{}
\theoremstyle{dotlessdr}
\newtheorem{example}{Exemplu}[section]
\newtheorem{definition}{Defini\c{t}ie}[section]
\newtheorem{remark}{Observa\c{t}ie}[section]




\begin{document}

\begin{center}
  {\large\textbf{Seminar 13 \\ Recapitulare}}
\end{center}

\textbf{MODEL 1}

1. Verificați dacă funcția $f(x, y) = (x + y)^{-\frac{1}{2}}$ este armonică.

\vspace{1cm}

2. Fie funcția $f(x, y) = x^3y + e^{x^2 + y^2}$. Calculați, folosind definiția,
derivatele parțiale în $(-1, 1)$.

\vspace{1cm}

3. Fie funcția $f(x, y) = e^{x + y}$. Să se scrie polinomul Taylor de gradul 5
asociat funcției $f$ în punctele $(0, 0)$ și $(1, -1)$.

\vspace{1cm}

4. Să se determine valorile extreme pentru funcția $f$, definită pe domeniul $D$,
unde:
\[
  f(x, y) = x^4 + y^3 - 4x^3 - 3y^2 + 3y, \quad
  D = \{(x, y) \in \RR^2 \mid x^2 + y^2 \leq 4 \}.
\]

\vspace{1cm}

5. Să se calculeze dreapta de regresie care mediază între punctele
$A(-1, 0), B(2, 3), C(3, 5)$.

\vspace{2cm}

\textbf{MODEL 2}

1. Fie funcția:
\[
  f(x, y) = \begin{cases}
    \frac{xy^3}{x^2 + y^2}, & (x, y) \neq (0, 0) \\
    0, & (x, y) = (0, 0)
  \end{cases}
\]

Arătați că funcția este de clasă $\kal{C}^1$ pe $\RR^2$.

\vspace{1cm}

2. Fie $f : \RR^3 - \{(0, 0, 0) \}, f(x,y,z) = \ln(x^2 + y^2 + z^2)$. Să se
calculeze $\Delta f$.

\vspace{1cm}

3. Să se calculeze aproximarea liniară în jurul originii a funcției:
\[
  f(x, y, z) = \sqrt{\frac{x + 1}{(y + 1)(z + 1)}}.
\]

\vspace{1cm}

4. Să se determine punctele cele mai apropiate de origine care se găsesc pe
suprafața de ecuație:
\[
  4x^2 + y^2 + z^2 - 8x - 4y + 4 = 0.
\]

\vspace{1cm}

5. Să se găsească dreapta de regresie care mediază între punctele $A(3, 7), B(-1, 1)$ și
$C(0, 3)$.

\newpage

\textbf{MODEL 3}

1. Fie funcția:
\[
  f(x, y, z) = g(xy, xyz + \frac{x^2}{2} + \frac{x^4}{4}).
\]
Să se arate că ea satisface ecuația:
\[
  -xy \frac{\partial f}{\partial x} + y^2 \frac{\partial f}{\partial y} + %
  x(1 + x^2) \frac{\partial f}{\partial z} = 0.
\]

\vspace{1cm}

2. Să se calculeze diferențiala de ordinul întîi și de ordinul al doilea pentru
funcția $f(x, y) = \arctan \frac{y}{x}$, definită pe domeniul maxim de definiție.

\vspace{1cm}

3. Fie funcția $f(x, y) = e^x \sqrt[4]{y}$.
\begin{enumerate}[(a)]
\item Să se determine aproximarea liniară și pătratică a funcției;
\item Să se calculeze aproximativ $e^{-0,2}\sqrt[4]{1,02}$, folosind polinomul
  Taylor de gradul al doilea.
\end{enumerate}

\vspace{1cm}

4. Să se determine valoarea maximă a produsului $xy$, dacă $x$ și $y$ sînt coordonatele
unui punct de pe elipsa de ecuație:
\[
  6x^2 + 8y^2 - 12 = 0.
\]

\vspace{1cm}

5. Să se determine dreapta de regresie care mediază între punctele $A(3, 1), B(5, 2)$ și
$C(-2, -1)$.

\end{document}
